\section{Introduction}
\label{sec:intro}

Buildings take roughly 40~per~cent of final energy in the European Union, most of it
burned for space heating and hot water, which places them at the centre of European
climate policy \citep{IEA2022buildings,EuropeanCommission2020GreenDeal}. Hundreds of
billions of euros of that heating plant will be replaced by 2050, and whether that capital
electrifies the heat or reuses the gas network for hydrogen is a choice slow to unwind. A
clean option exists, but the difficulty lies in delivery, because heat reaches
hundreds of millions of dwellings of differing vintage and fabric whose appliances are
replaced only as they fail \citep{Azevedo2022,Chatterjee2024Navigating}.

European policy has converged on heat pumps as the principal route, and the case is strong.
A coefficient of performance (COP) of two to four draws less final energy than any combustion
appliance, and the machine decarbonises as the grid does
\citep{IEA2022heatpumps,Abbasi2021}. However, deployment is slowed by frictions that fall
unevenly across the stock, among them retrofitting poorly insulated dwellings,
distribution-grid reinforcement in dense areas, and performance that degrades in the
coldest weather \citep{Lingard2020Residential,Gibb2023Coming}. We price that reinforcement
inside the levelised cost of delivered heat, and the same frictions open pockets for a
gas-based alternative.

Hydrogen is the alternative most often advanced, and the most contested, because a
low-carbon gas could in principle reuse the existing distribution network and the
appliance habit of a gas-heated home \citep{Dodds2015,Bothe2021}. REPowerEU raised the EU's 2030 ambition to 20~Mt of renewable and imported supply, with
buildings among the candidate demand centres
\citep{REpowerEU2022,EUCommissionSWD2022}, and the 2022 gas crisis sharpened the
appetite for a storable domestic fuel \citep{Goldthau2023Eu}. However, a substantial
sceptical literature answers that hydrogen for space heating wastes primary energy,
competes for scarce electrolytic supply, and understates the cost of network conversion
\citep{Rosenow2022PipeDream,Rosenow2024MetaReview,Cebon2022,Badakhsh2023Caveats,ICCT2023h2heating}.

The quantitative basis for settling the dispute is thinner than the debate implies, and
the gaps fall into three groups. The first is spatial resolution. The continental
assessments resolve the building stock in detail but at national resolution, so they
cannot see the within-country demand gradient or the storage geology that decide where
hydrogen might compete \citep{Billerbeck2023Race,Brown2018Synergies}. Sub-national
European work exists for the hydrogen network itself
\citep{Neumann2023H2Network,Kountouris2024Unified}, spatially explicit studies are mostly
single-country \citep{Maestre2022,Huckebrink2022,Hoseinpoori2022}, and the pan-European
mapping resolves heat demand without a dispatch test
\citep{Mller2019Heat,Persson2019Heat}. The second is the share itself, which studies
commonly fix by assumption \citep{Weidner2023Planetary,Quarton2019Curious} and price on
annual levelised cost, which misstates how a dispatchable fuel competes
\citep{Joskow2008capacity,Quarton2019Value}. The third is a severed link between heat and
power, and coupled system studies do model it
\citep{Zeyen2021Mitigating,Mellot2024Mitigating}. However, they report a system optimum
rather than the standalone capital-recovery position of a hydrogen peaker market by
market, which is the quantity a private investor and a capacity mechanism both need.

Against those strands, this paper adds a quantity each leaves out. The review literature
concludes against hydrogen for domestic heat
\citep{Rosenow2022PipeDream,Rosenow2024MetaReview} without reporting where the conclusion
is close, so we reproduce the ordering in 22 of 29 markets, at an all-29 median of
\euro{}15.4/MWh in the heat pump's favour, and price the seven exceptions individually.
The sub-national work returns a network plan rather than a plant's balance sheet
\citep{Neumann2023H2Network,Kountouris2024Unified}, so we return a peaker's standalone
capital recovery market by market. The firm-capacity literature establishes that
low-carbon systems need clean firm capacity
\citep{Sepulveda2018Firm,Bistline2020ValueOfTechnology} without pricing it when the load
is building heat, so we put that at \euro{}31 to 63/kW-yr.

First, where other work imposes a hydrogen share and stops, we impose transparent adoption pathways and then test them against an independent dispatch screen, market by market, exposing the cost and geology regimes a continental average would bury. Second, we recast the question from an energy-share contest into one about capacity and capital, screening the linked tests on one dispatch basis and testing capital recovery on the plant that wins the power peak. To deliver both we assemble a bottom-up techno-economic model of residential heating across the EU-27, the United Kingdom and Switzerland from 2025 to 2050. It reconstructs useful heat demand from NUTS3 building stock and aggregates it to the country before any cost test, benchmarking to within $-0.8$~per~cent against an independent regional database \citep{MilojevicDupont2023Eubucco,Loga2016Tabula}, and propagates that demand through eight technologies and four scenarios in a 200-draw Monte Carlo.

On the levelised comparison we price the two carriers on a symmetric distribution basis, charging
the heat pump for the low- and medium-voltage reinforcement its winter peak requires and
the hydrogen boiler for its own last-mile network and seasonal store, and we size
equipment capital to peak heat demand. Published comparisons differ in which of these they charge, so we sweep each. Adopting them narrows the heat pump's margin without reversing the median ordering. However, corners of the sweep do turn it.

On that accounting the best 2050 heat pump undercuts the gas boiler in 24 of 29 countries
and the hydrogen boiler in 22. Cheap storage widens those margins without selecting them, since five of the seven still hold when repriced at the non-cavern rate \citep{Caglayan2020caverns}. Repriced on the short-run marginal cost (SRMC) a system operator
uses, hydrogen wins the cold-snap peak in a minority of markets that grows with policy ambition. However, no market recovers that peaker's capital, about 9~per~cent in the central case, the classic missing-money problem.

These findings answer one research question: under what conditions, and where, is
decarbonised hydrogen competitive against heat-pump electrification in the buildings of the
EU-27, the United Kingdom and Switzerland by 2050? Switzerland, a high-cost, hydro-rich and cavern-poor system reliant on imports,
provides a compact standalone case \citep{Gabrielli2020,IEA2023switzerland}.

The remainder is organised as follows. Section~\ref{sec:method} sets out the model and its
weather series, Section~\ref{sec:results} the three tests and two cross-checks, and
Section~\ref{sec:conclusions} the conclusions and policy implications.
